\corrigesChapitre{12}{Statistiques à deux variables}
\begin{corrige}{1}{Point moyen}
\begin{questions}
\item $\bar x=(1+2+3)/3=2$.
\item $\bar y=(3+4+8)/3=5$.
\item $G(2,5)$ ; il se calcule séparément sur les deux coordonnées.
\end{questions}
\end{corrige}
\begin{corrige}{2}{Utiliser une droite}
\begin{questions}
\item $2{,}5\times6+7=22$.
\item $2{,}5\times20+7=57$.
\item La première est une interpolation ; la seconde une extrapolation éloignée, plus incertaine et à justifier.
\end{questions}
\end{corrige}
\begin{corrige}{3}{Résidu}
\begin{questions}
\item $3(2)+1=7$.
\item $8-7=1$.
\item $1^2=1$.
\end{questions}
\end{corrige}
\begin{corrige}{4}{Lire une corrélation}
\begin{questions}
\item Une forte association affine décroissante dans les données étudiées.
\item Il ne prouve ni une causalité ni une validité du modèle hors du domaine observé.
\item Par exemple, des points symétriques d’une parabole peuvent avoir une covariance nulle tout en suivant une relation déterministe.
\end{questions}
\end{corrige}
\begin{corrige}{5}{Calculer une régression}
\begin{questions}
\item $\bar x=2{,}5$, $\bar y=4$, $V_x=1{,}25$, $\operatorname{Cov}=1{,}75$.
\item $a=1{,}75/1{,}25=1{,}4$, $b=4-1{,}4(2{,}5)=0{,}5$.
\item $r=1{,}75/\sqrt{1{,}25\times2{,}5}\simeq0{,}989949$.
\end{questions}
\end{corrige}
\begin{corrige}{6}{Exponentielle par transformation}
\begin{questions}
\item $z=\ln5+x\ln2$.
\item Pente $\ln2$ ; ordonnée $\ln5$.
\item $y=\mathrm e^{\ln5}\mathrm e^{x\ln2}=5\times2^x$.
\end{questions}
\end{corrige}
\begin{corrige}{7}{Interpolation et seuil}
\begin{questions}
\item $4\mathrm e^{0{,}9}\simeq9{,}838$.
\item $\mathrm e^{0{,}3t}=5$, donc $t=\ln5/0{,}3\simeq5{,}365$.
\item Cette date dépasse le domaine observé ; c’est une extrapolation, à considérer avec prudence.
\end{questions}
\end{corrige}
\begin{corrige}{8}{Moindres carrés}
\begin{questions}
\item $S_1=1+1+0=2$ et $S_2=4$.
\item $D_1$, car $2<4$.
\item Non : on n’a comparé que deux candidates, sans minimiser sur l’ensemble des droites.
\end{questions}
\end{corrige}
\begin{corrige}{9}{Corrélation nulle et parabole}
\begin{questions}
\item $G(0,2/3)$.
\item $[(-1)(1-2/3)+0+(1)(1-2/3)]/3=0$.
\item $y=x^2$ pour les trois points. Le coefficient de corrélation affine nul n’exclut pas une liaison non linéaire.
\end{questions}
\end{corrige}
\begin{corrige}{10}{Une variable confondante}
\begin{questions}
\item La température ou la saison peut influencer les deux.
\item Non : la corrélation seule ne suffit pas à attribuer une causalité directe.
\item Le protocole, la chronologie, les facteurs communs et les mécanismes possibles ; comparer des situations où ces facteurs sont contrôlés.
\end{questions}
\end{corrige}
\begin{corrige}{11}{Changer les unités}
\begin{questions}
\item $x=u/60$, donc $y=u+5$.
\item Non : les nombres de l’abscisse changent avec l’unité, mais la situation et la distance prédite restent les mêmes.
\item Non : multiplier une variable par une constante positive laisse $r$ inchangé ; la pente numérique, elle, change.
\end{questions}
\end{corrige}
\begin{corrige}{12}{Pourquoi la droite passe par $G$}
\begin{questions}
\item Avec les variables centrées, $S(a,b)=\sum(Y_i-aX_i)^2+n(\bar y-a\bar x-b)^2$, car les termes croisés se compensent.
\item Le carré est minimal en zéro : $b=\bar y-a\bar x$, donc $a\bar x+b=\bar y$.
\item $V_x=0$ ; la formule de la pente divise par zéro et une pente unique n’est pas identifiable à partir d’un seul niveau de $x$.
\end{questions}
\end{corrige}
